\emph{સાબિતી-તંત્રો} વાક્યો વચ્ચેની ફલિતતા અને સુસંગતતાને ઓળખવાની શુદ્ધ વાક્યરચનાત્મક રીતો આપે છે. \emph{સિક્વન્ટ કલન} આવું એક તંત્ર છે. તેમાં \emph{નિષ્પત્તિ} સિક્વન્ટોનું વૃક્ષ છે. સિક્વન્ટ $\Gamma\Sequent\Delta$માં સૂત્રોની બે શ્રેણી $\Sequent$થી અલગ પડે છે. વૃક્ષમાં મોખરે આવતા સિક્વન્ટો $!A\Sequent !A$ આકારના \emph{આરંભિક સિક્વન્ટો} છે. નિષ્પત્તિ સાચી હોય તે માટે બીજા બધા સિક્વન્ટોનું સમર્થન કોઈ \emph{અનુમાન-નિયમ}થી થવું જોઈએ. દરેક સંયોજક અને પરિમાણક માટે સિક્વન્ટની ડાબી અને જમણી બાજુ કામ કરતા નિયમોની જોડ હોય છે. દાખલા તરીકે, સિક્વન્ટ $\Gamma\Sequent\Delta,!A\lif !B$ને $\RightR{\lif}$ નિયમથી મેળવ્યો હોય તો તેનો \emph{આધાર સિક્વન્ટ} $!A,\Gamma\Sequent\Delta,!B$ હોવો જોઈએ. અમુક નિયમો સિક્વન્ટમાં સૂત્રોનો ક્રમ કે સંખ્યા બદલવા દે છે. દાખલા તરીકે, $\RightR{\Exchange}$ નિયમ જમણી બાજુનાં બે સૂત્રોની અદલાબદલી કરવા દે છે.

$\quad\Sequent !A$ની નિષ્પત્તિ હોય તો $!A$ \emph{પ્રમેય} છે, $\Proves !A$. $\Gamma_0$નું દરેક $!B$ એ $\Gamma$માં હોય અને $\Gamma_0\Sequent !A$ની નિષ્પત્તિ હોય તો $!A$ \emph{$\Gamma$માંથી નિષ્પન્ન કરી શકાય છે}, $\Gamma\Proves !A$. આવી $\Gamma_0$ માટે $\Gamma_0\Sequent\quad$ની નિષ્પત્તિ હોય તો $\Gamma$ \emph{અસંગત} છે, અન્યથા \emph{સુસંગત} છે. આ ખ્યાલો પરસ્પર જોડાયેલા છે. દાખલા તરીકે, $\Gamma\Proves !A$ ત્યારે અને માત્ર ત્યારે જ્યારે $\Gamma\cup\{\lnot !A\}$ અસંગત હોય. તેમનો અર્થવિચારી ખ્યાલો સાથે પણ સંબંધ છે. દાખલા તરીકે, $\Gamma\Proves !A$ હોય તો $\Gamma\Entails !A$. $\Gamma$માંથી નિષ્પન્ન કરી શકાય તે $\Gamma$માંથી ફલિત પણ થાય છે એવા ગુણધર્મને \emph{યથાર્થતા} કહે છે. \emph{યથાર્થતા પ્રમેય} નિષ્પત્તિની લંબાઈ પર આગમન વડે સાબિત થાય છે. દરેક અનુમાનના આધાર સિક્વન્ટો પ્રામાણ્ય ધરાવતા હોય ત્યારે નિષ્કર્ષ સિક્વન્ટ પણ પ્રામાણ્ય ધરાવે છે એમ બતાવાય છે.
